\chapter[4.1 Requisitos y especificaciones]{4.1 Análisis de requisitos y especificaciones técnicas}
\label{ch:requisitos}

Los requisitos se derivan de las restricciones físicas de la celda, los criterios de seguridad de la norma IEC~61511 y los objetivos de producción. El sistema de control prioriza en este orden: seguridad operativa (cero desbordamientos), estabilidad de la ocupación en torno al nivel objetivo, y eficiencia energética con máxima productividad.

\section{4.1.a Requisitos funcionales y no funcionales}

\subsection*{Requisitos funcionales (RF)}

\begin{longtable}{@{}P{0.11\linewidth}P{0.23\linewidth}P{0.55\linewidth}@{}}
\caption{Requisitos funcionales del sistema de control inteligente R2ET.}
\label{tab:req_func}\\
\toprule
ID & Requisito & Descripción \\ \midrule \endfirsthead
\toprule ID & Requisito & Descripción \\ \midrule \endhead
RF-01 & Regulación de nivel & Mantener $n_i(k) \approx n^* = 1.5$~piezas en ambas esteras ajustando $u_i(k)$.\\
RF-02 & No desbordamiento & $n_i(k) \leq 3$~piezas en todo momento. Incumplimiento implica parada.\\
RF-03 & No vaciado & $n_i(k) \geq 0.2$~piezas para garantizar flujo continuo de salida.\\
RF-04 & Coordinación del robot & Ajustar $r_i(k)$ en cada ciclo según el desequilibrio entre esteras.\\
RF-05 & Rechazo de perturbación & Compensar el cambio de peso sin superar la alarma de 2.7~piezas.\\
RF-06 & Control de admisión & Reducir $a(k)$ cuando $n_{max} > 2.4$~piezas.\\
RF-07 & Seguimiento de referencia & Error estacionario $<5\%$ desde $n_i(0)=1.25$.\\
RF-08 & Señalización & Activar alarma cuando $n_i > 2.7$ y registrar tiempo en zona de riesgo.\\
RF-09 & Servoregulación local & Regular $y(k)$ ante escalón y perturbación en la planta de Fase~1.\\
RF-10 & Adaptación automática & Sin resintonización manual para perturbaciones dentro del rango de diseño.\\
\bottomrule
\end{longtable}

\subsection*{Requisitos no funcionales (RNF)}

\begin{longtable}{@{}P{0.11\linewidth}P{0.23\linewidth}P{0.55\linewidth}@{}}
\caption{Requisitos no funcionales del sistema de control.}
\label{tab:req_nofunc}\\
\toprule
ID & Requisito & Descripción \\ \midrule \endfirsthead
\toprule ID & Requisito & Descripción \\ \midrule \endhead
RNF-01 & Precisión & IAE (Fase~2) $<60$ en 160~muestras; error estacionario $<2\%$.\\
RNF-02 & Adaptabilidad & Control efectivo ante $\pm20\%$ de variación en caudal de entrada.\\
RNF-03 & Robustez & Cero violaciones en condiciones nominales y en escenario severo ($+20\%$~demanda, $-15\%$~salida).\\
RNF-04 & Seguridad & El supervisor de capacidad es independiente del controlador primario y siempre activo.\\
RNF-05 & Escalabilidad & La arquitectura supervisora debe admitir $N>2$~esteras con cambios localizados.\\
RNF-06 & Mantenibilidad & Cada método (PID, Fuzzy, DL, QL) es intercambiable sin tocar la capa supervisora.\\
RNF-07 & HMI & Panel de operador con $n_1$, $n_2$, alarmas y estado del robot en tiempo real.\\
RNF-08 & Integración & Compatible con OPC~UA para conexión con SCADA/MES de nivel superior.\\
RNF-09 & Tiempo de ciclo & Acción de control calculada dentro del ciclo de la celda R2ET ($1$~s).\\
RNF-10 & Trazabilidad & Registro temporal de todas las variables para análisis post-mortem.\\
\bottomrule
\end{longtable}

\section{4.1.b Especificaciones técnicas del sistema de control}

\subsection*{Variables de proceso y control}

\begin{longtable}{@{}P{0.18\linewidth}P{0.13\linewidth}P{0.13\linewidth}P{0.44\linewidth}@{}}
\caption{Variables de proceso y de control del sistema R2ET.}
\label{tab:var_proceso}\\
\toprule
Variable & Rango & Unidad & Descripción \\ \midrule \endfirsthead
\toprule Variable & Rango & Unidad & Descripción \\ \midrule \endhead
\multicolumn{4}{@{}l}{\textit{Variables de proceso (medidas)}} \\ \midrule
$n_1(k)$ & $[0,\;3]$ & piezas & Nivel de llenado estera E1\\
$n_2(k)$ & $[0,\;3]$ & piezas & Nivel de llenado estera E2\\
$y(k)$ & $[0,\;1.2]$ & norm. & Salida de la planta servo (Fase~1)\\
$\lambda(k)$ & $[0.42,\;0.77]$ & p/muestra & Caudal de entrada del robot\\
$d(k)$ & $\{0,\;1\}$ & binario & Señal de perturbación (cambio de peso)\\
\midrule
\multicolumn{4}{@{}l}{\textit{Variables de control (actuadores)}} \\ \midrule
$u_1(k)$ & $[0.05,\;1.0]$ & norm. & Velocidad de descarga E1\\
$u_2(k)$ & $[0.05,\;1.0]$ & norm. & Velocidad de descarga E2\\
$r_1(k)$ & $[0.20,\;0.80]$ & fracción & Fracción del robot hacia E1 ($r_2=1-r_1$)\\
$a(k)$ & $[0.72,\;1.0]$ & fracción & Factor de admisión de nuevas piezas\\
\bottomrule
\end{longtable}

\subsection*{Parámetros operativos}

\begin{longtable}{@{}P{0.30\linewidth}P{0.15\linewidth}P{0.43\linewidth}@{}}
\caption{Parámetros operativos y su justificación.}
\label{tab:params_op}\\
\toprule
Parámetro & Valor & Justificación \\ \midrule \endfirsthead
\toprule Parámetro & Valor & Justificación \\ \midrule \endhead
Nivel objetivo $n^*$ & 1.5~piezas & 50\% de la capacidad: centra la operación con margen en ambos sentidos.\\
Capacidad $N_{max}$ & 3~piezas & Restricción física de la longitud de la estera.\\
Velocidad nominal $u_{ss}$ & 0.457~(norm.) & Equilibrio entrada-salida en condición de caudal base.\\
Umbral de riesgo & 2.4~piezas & 80\% de capacidad: inicio del refuerzo supervisorio.\\
Umbral de alarma & 2.7~piezas & 90\% de capacidad: actuación de emergencia.\\
Pérdida E1 por atasco & 35\% & Reducción de descarga con piezas pesadas en E1.\\
Pérdida E2 por atasco & 17.5\% & Menor impacto por posición diferente en E2.\\
Período de muestreo & 1~muestra & Ciclo de control digital.\\
\bottomrule
\end{longtable}

\subsection*{Restricciones técnicas}

Las restricciones se ordenan por dureza:

\begin{enumerate}
  \item Capacidad (dura): $0 \leq n_i(k) \leq 3$. Ninguna violación es admisible.
  \item Actuador: $0.05 \leq u_i(k) \leq 1.0$. El mínimo evita la parada total de la estera.
  \item Reparto: $0.20 \leq r_1(k) \leq 0.80$. Ambas esteras reciben siempre algo de flujo.
  \item Admisión en emergencia: $a(k) \geq 0.72$. No se corta completamente la entrada.
  \item Tiempo de ciclo: la acción de control debe calcularse en $T_s=1$~muestra.
\end{enumerate}

\subsection*{Variables de monitoreo e indicadores de desempeño objetivo}

\begin{longtable}{@{}P{0.30\linewidth}P{0.20\linewidth}P{0.38\linewidth}@{}}
\caption{Objetivos cuantitativos de desempeño.}
\label{tab:obj_desempeno}\\
\toprule
Indicador & Objetivo & Criterio de aceptación \\ \midrule \endfirsthead
\toprule Indicador & Objetivo & Criterio \\ \midrule \endhead
Violaciones de capacidad & 0 & Condición absoluta.\\
Tiempo en alarma ($n_i>2.7$) & 0~muestras & Deseable; tolerable $\leq 2$.\\
IAE Fase~2 (160~muestras) & $<60$ & Suma de errores de ambas esteras.\\
IAE durante atasco & $<25$ & Solo en $k=95$--$118$.\\
Margen de seguridad & $>0.3$~piezas & Distancia mínima a la capacidad.\\
Recuperación post-atasco & $<15$~muestras & Para salir de la zona de riesgo.\\
Productividad acumulada & $>60$~piezas & Caudal total de salida en 160~muestras.\\
Desbalance medio & $<0.5$~piezas & Media de $|n_1(k)-n_2(k)|$.\\
\bottomrule
\end{longtable}
